\documentclass[11pt]{article}
\newcommand{\ListaTitulo}{Gabarito 02: Tendência central, quantis e boxplot}
% Preâmbulo compartilhado: MATF14, Listas de Exercícios 2026
% Usado por ListaNN.tex e GabaritoNN.tex via % Preâmbulo compartilhado: MATF14, Listas de Exercícios 2026
% Usado por ListaNN.tex e GabaritoNN.tex via % Preâmbulo compartilhado: MATF14, Listas de Exercícios 2026
% Usado por ListaNN.tex e GabaritoNN.tex via \input{preamble}

\usepackage[utf8]{inputenc}
\usepackage[T1]{fontenc}
\usepackage[brazil]{babel}
\usepackage{fullpage}
\usepackage{amsmath,amssymb}
\usepackage{enumitem}
\usepackage{xcolor}
\usepackage{fancyhdr}
\usepackage{titlesec}
\usepackage{listings}
\usepackage{hyperref}
\usepackage{booktabs}
\usepackage{graphicx}
\usepackage{tikz}

\definecolor{matf14azul}{HTML}{0D3B54}
\definecolor{matf14dourado}{HTML}{C9971E}

\titleformat{\section}{\normalfont\Large\bfseries\color{matf14azul}}{\thesection}{1em}{}
\titleformat{\subsection}{\normalfont\large\bfseries\color{matf14azul}}{\thesubsection}{1em}{}

\providecommand{\ListaTitulo}{Lista de Exercícios}

\setlength{\headheight}{14pt}
\pagestyle{fancy}
\fancyhf{}
\lhead{\small MATF14: Estatística Econômica I}
\rhead{\small \ListaTitulo}
\cfoot{\thepage}
\renewcommand{\headrulewidth}{0.4pt}

\lstset{
  basicstyle=\ttfamily\small,
  breaklines=true,
  frame=single,
  columns=fullflexible,
  backgroundcolor=\color{gray!8},
  commentstyle=\color{matf14dourado},
  keywordstyle=\color{matf14azul}\bfseries,
  language=R
}

\newcounter{questaocounter}
\newenvironment{questao}[1]{%
  \stepcounter{questaocounter}%
  \bigskip\noindent\textbf{Questão #1.}\ %
}{\par}

\newcommand{\cabecalho}[2]{%
  \begin{center}
    {\Large\bfseries\color{matf14azul} #1}\\[0.3em]
    {\large #2}
  \end{center}
  \vspace{0.5em}
  \noindent\rule{\linewidth}{0.6pt}
  \vspace{0.5em}
}

\newenvironment{solucao}{%
  \par\smallskip\noindent\textit{\color{matf14dourado!80!black}Solução.}\ %
}{\hfill$\blacksquare$\par}


\usepackage[utf8]{inputenc}
\usepackage[T1]{fontenc}
\usepackage[brazil]{babel}
\usepackage{fullpage}
\usepackage{amsmath,amssymb}
\usepackage{enumitem}
\usepackage{xcolor}
\usepackage{fancyhdr}
\usepackage{titlesec}
\usepackage{listings}
\usepackage{hyperref}
\usepackage{booktabs}
\usepackage{graphicx}
\usepackage{tikz}

\definecolor{matf14azul}{HTML}{0D3B54}
\definecolor{matf14dourado}{HTML}{C9971E}

\titleformat{\section}{\normalfont\Large\bfseries\color{matf14azul}}{\thesection}{1em}{}
\titleformat{\subsection}{\normalfont\large\bfseries\color{matf14azul}}{\thesubsection}{1em}{}

\providecommand{\ListaTitulo}{Lista de Exercícios}

\setlength{\headheight}{14pt}
\pagestyle{fancy}
\fancyhf{}
\lhead{\small MATF14: Estatística Econômica I}
\rhead{\small \ListaTitulo}
\cfoot{\thepage}
\renewcommand{\headrulewidth}{0.4pt}

\lstset{
  basicstyle=\ttfamily\small,
  breaklines=true,
  frame=single,
  columns=fullflexible,
  backgroundcolor=\color{gray!8},
  commentstyle=\color{matf14dourado},
  keywordstyle=\color{matf14azul}\bfseries,
  language=R
}

\newcounter{questaocounter}
\newenvironment{questao}[1]{%
  \stepcounter{questaocounter}%
  \bigskip\noindent\textbf{Questão #1.}\ %
}{\par}

\newcommand{\cabecalho}[2]{%
  \begin{center}
    {\Large\bfseries\color{matf14azul} #1}\\[0.3em]
    {\large #2}
  \end{center}
  \vspace{0.5em}
  \noindent\rule{\linewidth}{0.6pt}
  \vspace{0.5em}
}

\newenvironment{solucao}{%
  \par\smallskip\noindent\textit{\color{matf14dourado!80!black}Solução.}\ %
}{\hfill$\blacksquare$\par}


\usepackage[utf8]{inputenc}
\usepackage[T1]{fontenc}
\usepackage[brazil]{babel}
\usepackage{fullpage}
\usepackage{amsmath,amssymb}
\usepackage{enumitem}
\usepackage{xcolor}
\usepackage{fancyhdr}
\usepackage{titlesec}
\usepackage{listings}
\usepackage{hyperref}
\usepackage{booktabs}
\usepackage{graphicx}
\usepackage{tikz}

\definecolor{matf14azul}{HTML}{0D3B54}
\definecolor{matf14dourado}{HTML}{C9971E}

\titleformat{\section}{\normalfont\Large\bfseries\color{matf14azul}}{\thesection}{1em}{}
\titleformat{\subsection}{\normalfont\large\bfseries\color{matf14azul}}{\thesubsection}{1em}{}

\providecommand{\ListaTitulo}{Lista de Exercícios}

\setlength{\headheight}{14pt}
\pagestyle{fancy}
\fancyhf{}
\lhead{\small MATF14: Estatística Econômica I}
\rhead{\small \ListaTitulo}
\cfoot{\thepage}
\renewcommand{\headrulewidth}{0.4pt}

\lstset{
  basicstyle=\ttfamily\small,
  breaklines=true,
  frame=single,
  columns=fullflexible,
  backgroundcolor=\color{gray!8},
  commentstyle=\color{matf14dourado},
  keywordstyle=\color{matf14azul}\bfseries,
  language=R
}

\newcounter{questaocounter}
\newenvironment{questao}[1]{%
  \stepcounter{questaocounter}%
  \bigskip\noindent\textbf{Questão #1.}\ %
}{\par}

\newcommand{\cabecalho}[2]{%
  \begin{center}
    {\Large\bfseries\color{matf14azul} #1}\\[0.3em]
    {\large #2}
  \end{center}
  \vspace{0.5em}
  \noindent\rule{\linewidth}{0.6pt}
  \vspace{0.5em}
}

\newenvironment{solucao}{%
  \par\smallskip\noindent\textit{\color{matf14dourado!80!black}Solução.}\ %
}{\hfill$\blacksquare$\par}


\begin{document}

\cabecalho{Gabarito 02}{Medidas de tendência central, quantis, IQR e boxplot (Cap. 1, Aulas 04--05)}

\begin{questao}{1} \textbf{(Cálculo direto: salários de uma pequena empresa)}
\begin{solucao}
\begin{enumerate}[label=(\alph*)]
  \item Soma $=3+3+3{,}5+4+4+4{,}5+5+6+22=55$, $n=9$, logo $\bar x = 55/9 \approx 6{,}11$.
    Ordenados, o valor central (5º de 9) é $4$: mediana $=4$. Moda: $3$ aparece 2 vezes, $4$
    aparece 2 vezes, bimodal (3 e 4).
  \item A mediana (R\$ 4 mil) representa melhor o "salário típico": a média (R\$ 6,11 mil) está
    inflada pelo único salário muito alto (R\$ 22 mil, um outlier claro em relação ao resto),
    enquanto a mediana não é afetada por esse valor.
  \item A média mudaria mais (nova soma $=41$, $\bar x = 41/9\approx 4{,}56$, uma queda grande). A
    mediana não mudaria nada: o valor na posição central continua sendo $4$ independentemente do
    valor do maior salário, contanto que ele continue sendo o maior.
\end{enumerate}
\end{solucao}
\end{questao}

\begin{questao}{2} \textbf{(Assimetria e ordenação das medidas)}
\begin{solucao}
\begin{enumerate}[label=(\alph*)]
  \item Como média (340) $>$ mediana (210), a distribuição é assimétrica à \textbf{direita}: existe uma
    cauda de famílias com patrimônio muito alto, puxando a média para cima, enquanto a mediana
    continua próxima do "centro" da maioria das famílias.
  \item Um histograma com um pico concentrado em valores relativamente baixos/moderados e uma
    cauda longa e fina se estendendo para a direita, em direção a valores de patrimônio muito
    altos (poucas famílias muito ricas).
  \item A mediana (R\$ 210 mil), por não ser distorcida pela cauda de patrimônios muito altos, ela
    descreve melhor a experiência da família "do meio" da distribuição, que é geralmente o que se
    quer comunicar como "típico".
\end{enumerate}
\end{solucao}
\end{questao}

\begin{questao}{3} \textbf{(Quantis: tempo de espera em atendimento bancário)}
\begin{solucao}
Dados ordenados, $n=12$: $2,3,3,5,6,7,8,9,10,12,15,30$.
\begin{enumerate}[label=(\alph*)]
  \item Usando interpolação simples (posições $0{,}25(n+1)=3{,}25$ e $0{,}75(n+1)=9{,}75$):
    $Q_1 \approx 3 + 0{,}25\times(5-3) = 3{,}5$; $Q_2 = (7+8)/2 = 7{,}5$;
    $Q_3 \approx 10 + 0{,}75\times(12-10) = 11{,}5$. (Pequenas variações são aceitáveis dependendo
    do método de interpolação usado, o importante é a lógica de posição.)
  \item $IQR = Q_3 - Q_1 \approx 11{,}5 - 3{,}5 = 8{,}0$.
  \item Limite superior $= Q_3 + 1{,}5\times IQR \approx 11{,}5 + 1{,}5\times 8{,}0 = 23{,}5$.
    Como $30 > 23{,}5$, o valor 30 minutos \textbf{é} classificado como outlier pelo critério do
    boxplot.
\end{enumerate}
\end{solucao}
\end{questao}

\begin{questao}{4} \textbf{(Leitura de boxplot: comparação entre grupos)}
\begin{solucao}
\begin{enumerate}[label=(\alph*)]
  \item A agência A, com caixa (IQR) mais estreita, tem processo mais consistente/previsível: a
    maioria das propostas leva um tempo parecido para ser aprovada. B tem maior variabilidade.
  \item Não necessariamente. A mediana menor de B indica que, na "propostas típica", B é um pouco
    mais rápida, mas os outliers mostram que uma fração não desprezível de propostas em B demora
    muito mais (acima de 20 dias), o que pode ser pior para o cliente médio em termos de risco de
    espera excessiva, mesmo com mediana menor.
  \item Para A: manter o processo, que é estável. Para B: investigar as causas dos casos extremos
    (outliers), provavelmente um subprocesso ou tipo de proposta específico que trava, mesmo que
    o "caso comum" seja rápido.
\end{enumerate}
\end{solucao}
\end{questao}

\begin{questao}{5} \textbf{(Dedução: robustez da mediana)}
\begin{solucao}
\begin{enumerate}[label=(\alph*)]
  \item Por definição, para $n$ ímpar a mediana é o valor na posição $(n+1)/2$ da lista
    \textbf{ordenada}. Se $x_n$ (o maior valor) é substituído por qualquer valor $x_n' > x_{n-1}$ (de
    modo que ele continue sendo o maior), a ordem relativa de todos os outros elementos, inclusive
    o da posição central $(n+1)/2 < n$, não se altera. Como a mediana depende apenas do valor
    naquela posição central, e esse valor não mudou, a mediana permanece a mesma.
  \item Exemplo com $n=5$: dados $1,2,3,4,5$. Média $=15/5=3$, mediana $=3$. Substituindo $x_5=5$
    por $x_5'=100$: novos dados $1,2,3,4,100$. Nova média $=110/5=22$ (mudou drasticamente),
    mas a mediana continua $3$ (não mudou).
  \item Como mostrado em (a) e (b), a mediana depende só da posição central dos dados ordenados,
    ignora completamente a magnitude dos valores fora dessa posição, por isso é chamada
    \textbf{robusta}. A média usa o valor exato de \emph{todos} os pontos na soma, então qualquer
    valor extremo (por maior que seja) afeta seu resultado proporcionalmente, ela \textbf{não} é
    robusta a outliers.
\end{enumerate}
\end{solucao}
\end{questao}

\end{document}
